\documentclass[sigconf]{acmart}

\usepackage{graphicx}
\usepackage{booktabs}
\usepackage{hyperref}
\usepackage{url}
\usepackage{caption}
\usepackage{subcaption}

\AtBeginDocument{%
  \providecommand\BibTeX{{%
    \normalfont B\kern-0.5em{\scshape i\kern-0.25em b}\kern-0.8em\TeX}}}

% Copyright, conference, and ISBN details removed — to be added before submission
\setcopyright{none}
\renewcommand\footnotetextcopyrightpermission[1]{} % suppress copyright footnote

% Suppress running headers (venue/date line)
\acmConference{}{}{}
\acmBooktitle{}
\acmPrice{}
\acmISBN{}
\acmDOI{}

% Suppress ACM Reference Format block and running headers
\settopmatter{printacmref=false, printfolios=true}

\begin{document}

\title{LLM-Based Intelligent Notification Composition: From Static Personalization to Context-Aware Persuasive Messaging}

\author{Nilesh Agrawal}
\email{nilesh.d.agrawal@gmail.com}
\orcid{0000-0000-0000-0000}
\affiliation{%
  \institution{Independent Researcher}
  \city{Seattle}
  \state{WA}
  \country{USA}
}

\renewcommand{\shortauthors}{}
\fancyhead{} % clear all header fields

\begin{abstract}
Push notifications remain among the most direct and time-sensitive channels through which digital platforms engage users, yet conventional approaches to notification copywriting---manual authoring and rule-based templating---do not scale to the billions of user interactions generated by modern platforms. This paper presents a systematic survey and architectural framework for LLM-based intelligent notification composition. Modern notification systems are highly effective at deciding \textit{who} to notify, \textit{when} to notify, and \textit{what} to recommend, but they remain weak at deciding \textit{how} to communicate. Most systems still rely on rigid, slot-filled templates that produce relevance without intelligence. We argue that the next evolution of notifications is not simply better personalization, but intelligent composition: using LLMs to transform static alerts into context-aware, persuasive, and adaptive messages. We examine how LLMs are being integrated into production notification pipelines across social media, food delivery, and e-commerce, focusing on the technical mechanisms that enable scalable and context-aware generation, including Retrieval-Augmented Generation (RAG), Parameter-Efficient Fine-Tuning (PEFT), pairwise reward modeling, and contextual bandits for send-time optimization.

A central contribution of this work is the explicit disentanglement of message generation from adjacent recommender system components such as audience targeting, content ranking, and delivery timing. We argue that observed gains in engagement are often incorrectly attributed to text generation alone, when in practice they arise from interaction effects across the full pipeline. Rather than proposing a novel algorithmic primitive, this paper contributes an integration-level architectural synthesis: a unified notification framework that incorporates budget-aware routing, grounded generation, candidate ranking, diversity controls, frequency capping, and online learning. We also critically assess current evaluation practice, highlighting the mismatch between short-term click metrics and long-term user health, the fragility of causal inference assumptions in networked platforms, and the ethical boundary between personalization and manipulation. Our analysis suggests that LLMs can deliver meaningful value in high-context, high-variance notification settings, but only when deployed with rigorous safeguards, careful experimental design, and explicit recognition that language generation is not always the primary bottleneck.
\end{abstract}

\begin{CCSXML}
<ccs2012>
   <concept>
       <concept_id>10002951.10003317.10003347.10003350</concept_id>
       <concept_desc>Information systems~Recommender systems</concept_desc>
       <concept_significance>500</concept_significance>
       </concept>
   <concept>
       <concept_id>10010147.10010178.10010179.10010182</concept_id>
       <concept_desc>Computing methodologies~Natural language generation</concept_desc>
       <concept_significance>500</concept_significance>
       </concept>
   <concept>
       <concept_id>10003120.10003130.10011762</concept_id>
       <concept_desc>Human-centered computing~Empirical studies in collaborative and social computing</concept_desc>
       <concept_significance>300</concept_significance>
       </concept>
 </ccs2012>
\end{CCSXML}

\ccsdesc[500]{Information systems~Recommender systems}
\ccsdesc[500]{Computing methodologies~Natural language generation}
\ccsdesc[300]{Human-centered computing~Empirical studies in collaborative and social computing}

\keywords{Large Language Models, Push Notifications, Personalization, Natural Language Generation, Recommender Systems, Causal Inference, Retrieval-Augmented Generation, Reward Modeling}

\maketitle

\section{Introduction}

In the contemporary digital economy, user attention is scarce, interruptible, and intensely contested. Consumer platforms invest heavily in acquisition, but long-term value depends just as much on retention, reactivation, and repeated high-value engagement. Push notifications have therefore become a central operational mechanism for re-engaging users outside the application surface. On large consumer platforms, daily notification volumes range from millions to billions, and small changes in notification quality can materially affect click-through rates (CTR), conversion outcomes, and the broader health of the platform ecosystem \cite{bie2026pushgen} [T1]. Poorly targeted or repetitive notifications do not merely underperform; they actively erode trust, accelerate opt-outs, and contribute to long-term channel degradation \cite{ohly2023effects} [T1].

Historically, platforms have relied on two broad strategies for notification composition: manual copywriting and template-based generation. Most production notification systems today remain template-driven. Even when personalization is present, it is often limited to slot filling, such as inserting a user’s preferred category, recent item, or location into a fixed message pattern. This makes notifications relevant in a narrow sense, but rarely intelligent in how they communicate. The central tension is that modern systems are effective at deciding \textit{who} to notify, \textit{when} to notify, and \textit{what} to recommend, but they are still weak at deciding \textit{how} to communicate the recommendation in an adaptive, context-rich, and persuasive way.

Large Language Models (LLMs) create a new opportunity: not merely to personalize notifications, but to make them context-aware, adaptive, and semantically persuasive \cite{vaswani2017attention}. Consider a user with a historical preference for pizza receiving a dinner-time notification at 6 PM. A conventional system may send a static message such as ``Order your favorite pizza now.'' An LLM-enabled system, by contrast, can incorporate temporal context, prior taste signals, and controlled novelty to generate a richer message such as: ``It’s the perfect time for dinner---your usual pizza spot is nearby, and their sweet-and-savory pineapple special might be worth trying tonight.'' In this setting, the value of the LLM is not in identifying the user or the category alone, but in producing a more intelligent, context-sensitive, and action-inducing framing of the recommendation. This capability can also support controlled exploration by introducing adjacent or novel options in a way that feels coherent and persuasive, rather than random.

Yet enthusiasm for LLM-based messaging has also created analytical confusion. In many production pipelines, the notification message is only one stage in a larger decision stack that includes audience selection, content ranking, delivery timing, budget allocation, and post-hoc candidate filtering. Reported gains in engagement are often attributed broadly to ``AI notifications'' or ``LLM personalization,'' even when the observed uplift may be driven primarily by upstream improvements in targeting or timing. This conflation obscures both scientific understanding and engineering decision-making. A system that generates brilliant copy for the wrong user, about the wrong item, at the wrong time will still fail. Conversely, a well-ranked, well-timed notification may perform strongly even with relatively simple language realization.

This paper addresses that gap by providing both a systematic survey and an architectural framework for LLM-based intelligent notification composition. We argue that the next evolution of notifications is not simply better personalization, but intelligent composition: using LLMs to transform static or slot-filled alerts into context-aware, persuasive, and adaptive messages that better align recommendations with user intent, real-time context, and controlled novelty. The aim is not merely to summarize examples, but to clarify where language generation fits within the broader notification stack, what technical patterns currently dominate practice, how those patterns should be evaluated, and under what conditions LLM generation is actually worth deploying.

\subsection{Contributions and Novelty Statement}

This paper makes four contributions:

\begin{enumerate}
    \item \textbf{Systematic Synthesis:} It provides a structured systematic survey of academic and industrial work on LLM-based notification systems, using a PRISMA-guided review process and an evidence-tiering scheme to distinguish between peer-reviewed results, industrial deployment reports, and contextual sources.
    \item \textbf{Architectural Disentanglement:} It introduces an architectural attribution perspective that explicitly separates the contribution of message generation from adjacent components such as targeting, ranking, and send-time optimization.
    \item \textbf{Unified Framework:} It proposes a unified end-to-end framework for notification composition that integrates grounded generation, style adaptation, reward-based candidate ranking, diversity and frequency controls, budget-aware routing, and online learning. The novelty lies in \textit{integration novelty}---the logic and operational sequencing across known components---rather than the invention of any single component.
    \item \textbf{Critical Evaluation Framework:} It advances a critical evaluation and governance perspective by highlighting shortcomings in current measurement practice, especially around causal inference, heterogeneous treatment effects, long-term user health, and the ethical boundary between personalization and manipulation.
\end{enumerate}

\section{Review Methodology}

To support a structured and transparent synthesis, this survey followed a systematic review methodology designed for a fast-moving applied field in which relevant evidence is distributed across both academic publication channels and industrial documentation.

\subsection{Search Strategy and PRISMA Flow}

We conducted targeted searches across arXiv, ACM Digital Library, IEEE Xplore, and ScienceDirect, along with engineering publications and official technical blogs from major technology companies. The review covered the period from \textbf{January 2018 through March 2026}, chosen to capture the transition from pre-LLM notification optimization methods to contemporary generative systems. Search terms combined model-oriented concepts (``LLM'', ``RAG'', ``PEFT'', ``LoRA'', ``reward modeling'') with application-oriented phrases (``push notification'', ``message composition'', ``send-time optimization'', ``CTR optimization'', ``contextual bandit'').

The screening process followed a formal PRISMA (Preferred Reporting Items for Systematic Reviews and Meta-Analyses) flow, as illustrated in Figure \ref{fig:prisma}.

\begin{figure}[h]
    \centering
    \includegraphics[width=\linewidth]{fig1_prisma.png}
    \caption{PRISMA flow diagram detailing the literature search, screening, and inclusion process.}
    \label{fig:prisma}
\end{figure}

\subsection{Evidence Tiering and Source Limitations}

Because industrial notification systems often evolve more quickly than formal academic publication cycles, excluding non-academic sources would omit some of the most operationally relevant systems. At the same time, engineering blogs and deployment writeups frequently lack the methodological transparency needed for strong causal claims. To address this, sources were organized into three evidence tiers, which are marked inline throughout the text:

\begin{itemize}
    \item \textbf{[T1] Tier 1 (High Rigor):} Peer-reviewed papers and formally published preprints (e.g., WSDM, CHI, RecSys) with clear methodology, baseline comparisons, and meaningful technical depth.
    \item \textbf{[T2] Tier 2 (Moderate Rigor):} Official engineering blogs and deployment reports from major technology companies describing production systems and online results. \textbf{Limitation:} These often omit experimental controls or statistical detail.
    \item \textbf{[T3] Tier 3 (Contextual):} Vendor documentation and secondary summaries used only for framing broader trends rather than establishing performance claims.
\end{itemize}

\subsection{Corpus Summary and Quantitative Synthesis}

Table \ref{tab:corpus} summarizes the key primary sources anchoring our architectural synthesis, classified by domain and evaluation method. 

\begin{table*}[t]
\caption{Key Primary Sources and Systems Reviewed}
\label{tab:corpus}
\begin{tabular}{p{2.5cm}p{2cm}p{4cm}p{3cm}p{3cm}c}
\toprule
\textbf{System / Paper} & \textbf{Domain} & \textbf{Primary Technology} & \textbf{Evaluation Method} & \textbf{Key Finding} & \textbf{Tier} \\
\midrule
\textbf{PushGen \cite{bie2026pushgen}} & Social Media & LLM (SFT) + Pairwise Reward Model & Online A/B + Offline Accuracy & +14.08\% CTR on replaced content & [T1] \\
\textbf{Instagram Diversity \cite{sun2025ranking}} & Social Media & Similarity Penalty Layer & Online A/B & Reduced volume, improved CTR via diversity & [T2] \\
\textbf{DoorDash GNN \cite{sinha2024doordash}} & Food Delivery & ID-GNN + Contextual Copy & Online A/B & +1.8\% relative engagement lift & [T2] \\
\textbf{E-Commerce LLM \cite{acm2025ecommerce}} & E-Commerce & LLM Content Optimization & Online A/B & +12.5\% CTR, +8.3\% CVR & [T1] \\
\textbf{Tu et al. \cite{tu2021personalized}} & Professional Network & Heterogeneous Treatment Effects & Online A/B + Causal Simulation & Personalized volume caps reduce churn & [T1] \\
\textbf{Yancey et al. \cite{yancey2020sleeping}} & EdTech & Contextual Bandit (Thompson Sampling) & Online A/B & Send-time optimization increases DAU & [T1] \\
\textbf{Luo et al. \cite{luo2024survey}} & General RS & Causal Inference Survey & Literature Review & Identifies severe SUTVA violations in RS & [T1] \\
\bottomrule
\end{tabular}
\end{table*}

\textbf{Quantitative Synthesis:} Across the corpus, reported engagement lifts from LLM integration vary significantly by domain and baseline quality. Studies replacing manual or poorly optimized static templates report large CTR gains ranging from +8.0\% to +14.5\% \cite{bie2026pushgen, acm2025ecommerce} [T1]. However, in highly mature systems where LLMs are competing against heavily optimized contextual slot-filling, the marginal relative lifts are typically much smaller, clustering in the +1.0\% to +2.5\% range \cite{sinha2024doordash} [T2]. This variance underscores that ``LLM uplift'' is not an absolute constant, but a function of the baseline system's prior sophistication.

\section{Disentangling LLM Contributions from Adjacent Systems}

A recurring analytical error in both industrial discourse and academic reporting is to treat notification performance as if it were primarily a function of the copy itself. In real systems, however, the notification message is the output of a multi-stage decision process. The system must choose whom to notify, what to notify them about, when to deliver the message, how to phrase it, and, if multiple candidates are generated, which candidate to send. LLMs affect only a subset of that stack.

\begin{table}[h]
\caption{Architectural Attribution Matrix}
\label{tab:attribution}
\begin{tabular}{p{2.5cm}p{5cm}}
\toprule
\textbf{System Component} & \textbf{Contribution to Engagement} \\
\midrule
\textbf{Audience Selection} & Ensures budget is spent on users with highest marginal return; prevents spamming dormant users. \\
\textbf{Content Ranking} & Ensures the underlying item is highly relevant to the user's interests. \\
\textbf{Send-Time Optimization} & Capitalizes on user's daily routines; minimizes disruption. \\
\textbf{Message Generation (LLMs)} & Adapts tone, length, and specific hooks to user persona; increases semantic diversity. \\
\textbf{Candidate Ranking} & Aligns stylistic generation with historical CTR data; filters out poor LLM outputs. \\
\bottomrule
\end{tabular}
\end{table}

This decomposition (Table \ref{tab:attribution}) clarifies both the promise and the limits of LLMs. LLMs are strongest when the system already knows the relevant content and approximate delivery conditions, but still needs to decide how best to frame the notification. Their marginal value lies in semantic diversity, contextual nuance, and stylistic adaptation. They do not solve the targeting problem, the relevance problem, or the temporal decision problem. If those upstream components are weak, better generation may have only marginal impact.

\section{Core Architectural Frameworks and Technical Tradeoffs}

Across the literature, several technical patterns recur in the generation and ranking stages of notification systems. To move from descriptive summary to analytical synthesis, Table \ref{tab:synthesis} categorizes these dominant architectural patterns by their strengths, weaknesses, and primary failure modes.

\begin{table*}[t]
\caption{Comparative Synthesis of Architectural Patterns}
\label{tab:synthesis}
\begin{tabular}{p{3cm}p{4cm}p{4cm}p{3cm}}
\toprule
\textbf{Architectural Pattern} & \textbf{Main Strength} & \textbf{Main Weakness} & \textbf{Key Failure Mode} \\
\midrule
\textbf{Template / Slot Filling} & Low latency, high control, easy compliance & Low diversity, weak personalization & Message fatigue, generic wording \\
\textbf{RAG-Grounded LLM} & Context-aware personalization, factual grounding & Retrieval latency, grounding failures & Faithfulness hallucination \\
\textbf{PEFT / LoRA} & Cheap specialization, hot-swappable voice control & Requires curation of style objectives & Overfitting to narrow tone \\
\textbf{Pairwise Reward Modeling} & Better relative ranking than pointwise CTR prediction & Exposure bias, delayed reward issues & Reward hacking \\
\textbf{Diversity Penalty} & Reduces novelty fatigue, improves long-term quality & Can suppress short-term CTR maximization & Under-delivery of relevant items \\
\textbf{Contextual Bandit STO} & Adapts to user routines, improves timing efficiency & Exploration cost, sparse feedback & Wrong-time delivery for cold-start users \\
\textbf{Budget-Aware Routing} & Aligns compute with expected value & Depends on accurate CLV estimation & Over-serving already-active users \\
\textbf{Counterfactual Logging} & Supports off-policy evaluation and robust updates & Operationally costly and complex & Biased logs, unstable feedback \\
\bottomrule
\end{tabular}
\end{table*}

Table \ref{tab:synthesis} highlights that no single component is sufficient to explain observed performance gains in deployed notification systems. The empirical record suggests that improvements emerge from interaction effects across ranking, timing, generation, and governance layers. We now examine the technical depth of four core components.

\subsection{Retrieval-Augmented Generation and Grounding}

Retrieval-Augmented Generation (RAG) is one of the most natural fits for notification systems because notifications are highly contextual and often fact-sensitive \cite{lewis2020retrieval} [T1]. Rather than depending solely on the model's internal parameters, RAG injects external information into the prompt at inference time, typically including user profile attributes, item metadata, recent behavioral context, and campaign constraints. 

However, RAG introduces nontrivial system design challenges. Retrieval granularity affects both relevance and latency. Small chunks may improve precision but omit necessary surrounding context, while larger chunks improve recall at the expense of noise and prompt length \cite{mixofgranularity2025} [T1]. Furthermore, retrieval success does not guarantee faithful generation. Even when the correct context is retrieved, the model may ignore it, blend it with parametric memory, or produce an output that is only loosely supported by the evidence---a phenomenon known as ``faithfulness hallucination'' \cite{benchmarking2025} [T1]. The implication is that RAG should be treated as a grounding mechanism, not a complete factuality solution. Strong systems require downstream verification, either through rule-based checks, lightweight classifiers, or post-generation validation against structured context.

\subsection{Parameter-Efficient Fine-Tuning and Style Control}

Notification systems often require control over tone, brevity, urgency, and brand voice. While prompting alone can induce some stylistic variation, stable control at scale often requires model adaptation. Full fine-tuning is expensive and operationally inflexible. Parameter-Efficient Fine-Tuning (PEFT) methods, especially Low-Rank Adaptation (LoRA) \cite{hu2022lora} [T1], provide a more practical alternative by freezing the pre-trained model weights and injecting trainable rank decomposition matrices ($W = W_0 + BA$, where $B$ and $A$ are low-rank) into the Transformer layers.

This is particularly useful in notification systems because different campaigns, user segments, or brands may demand different stylistic priors. The appeal of LoRA is not only training efficiency but deployment flexibility: adapters can be swapped or merged with the base model weights during deployment, resulting in zero inference latency overhead \cite{han2024parameter} [T1]. Still, style adaptation creates its own risks. Over-optimization toward a narrow tone may increase short-term reward while pushing the system toward formulaic or manipulative language.

\subsection{Pairwise Reward Modeling and Candidate Ranking}

A central challenge in notification generation is not merely producing candidate messages, but identifying which candidate is most likely to produce the desired downstream action. Pointwise prediction of click probability from a single candidate is often unstable because observed outcomes are confounded by item quality, creator popularity, time of day, and historical policy bias. Pairwise reward modeling offers a more robust alternative: instead of predicting an absolute score, the model learns to prefer one candidate over another \cite{bie2026pushgen} [T1].

Yet reward modeling is also vulnerable to exposure bias. Historical logs contain feedback only for messages that were actually shown, meaning the reward model inherits the distributional blind spots of the previous system. Delayed rewards make the problem worse: the user may click immediately but convert much later, or not at all \cite{thompson2024} [T1]. Addressing this requires \textit{counterfactual logging}---deploying exploratory logging policies (like $\epsilon$-greedy) to gather unbiased data for off-policy evaluation \cite{joachims2018deep} [T1]. Without such mechanisms, reward models risk learning convenience rather than genuine causal usefulness.

\subsection{Send-Time Optimization Through Contextual Bandits}

Although notification copy is often the most visible element of the system, timing can be equally or more important than wording. Send-time optimization (STO) determines whether a message is delivered when the user is receptive, distracted, asleep, or oversaturated. 

Contextual multi-armed bandits provide a strong formalism for this task. By conditioning on user state and continuously trading off exploration and exploitation, bandit-based timing policies can discover individualized temporal windows for engagement. The Duolingo system models STO as a ``sleeping recovering bandit'' using Thompson Sampling for exploration \cite{yancey2020sleeping} [T1]. The ``state'' comprises user context (time zone, historical activity), the ``action'' is the decision to send or defer, and the ``reward'' is the engagement metric. The key lesson is architectural: timing should not be treated as a minor delivery detail after generation. In many systems it is one of the dominant determinants of performance \cite{ho2018notifying} [T1].

\section{Domain-Specific Applications and Limits of Generalization}

Although the underlying technical components recur across domains, the meaning of a ``good'' notification varies substantially by industry. Differences in user intent, action latency, business objective, and tolerance for persuasion shape both architecture and evaluation. Figure \ref{fig:domain} provides a conceptual mapping of these differences.

\begin{figure*}[t]
    \centering
    \includegraphics[width=\textwidth]{fig3_domain.png}
    \caption{Domain Generalizability Matrix comparing high-value user definitions, primary metrics, key contextual signals, LLM roles, and key risks across Social Media, Food Delivery, and E-Commerce.}
    \label{fig:domain}
\end{figure*}

\subsection{Social Media}

In social media platforms, notifications often drive content discovery, return sessions, or creator engagement. High-value users may be defined in terms of DAU, content production, or network centrality rather than direct revenue. This makes diversity especially important: repetitive or overly formulaic notifications may produce short-term clicks but harm long-term engagement through novelty fatigue. The PushGen system \cite{bie2026pushgen} [T1] and Instagram's diversity penalty layer \cite{sun2025ranking} [T2] are highly representative of this domain's need to balance relevance with semantic variety.

\subsection{Food Delivery and Local Commerce}

Food delivery notifications are strongly shaped by spatiotemporal context. Weather, time of day, delivery radius, restaurant availability, and recent ordering behavior all interact in ways that can make language framing materially important. A lunch offer framed around convenience during rain may outperform a generic discount message even when the underlying restaurant recommendation is the same \cite{sinha2024doordash} [T2]. This domain illustrates the critical importance of grounding: unlike pure content discovery, food delivery notifications encode strict operational claims about availability and ETA.

\subsection{E-Commerce}

E-commerce notifications are typically closer to direct-response marketing. Here the primary objective is not always the click itself, but conversion, order value, or recovery of partially completed user intent. Abandoned-cart flows, price-drop alerts, and browse-retargeting campaigns all provide natural contexts in which multiple framings can be plausible and economically meaningful \cite{acm2025ecommerce} [T1]. This makes e-commerce one of the strongest candidates for LLM-assisted generation, but also one of the most ethically delicate, as linguistic flexibility can easily drift toward manufactured urgency or artificial scarcity.

\section{When LLM Text Generation Is Not the Binding Constraint}

A recurring risk in the current discourse on LLM-driven notification systems is the assumption that better text generation necessarily yields the largest marginal gain in engagement. In mature notification stacks, this is often false. The dominant constraint may instead lie upstream in audience selection, content ranking, or send-time optimization. If the system selects the wrong user, promotes the wrong item, or delivers at the wrong time, even highly persuasive text will not recover performance.

This observation implies that LLM generation should be treated as a conditional optimization layer rather than a universal default. Its value is highest when the notification task has three properties: first, the promoted content admits multiple plausible framings; second, user response is sensitive to linguistic nuance rather than only item relevance; and third, the system has sufficient user context to support grounded personalization. These conditions are often satisfied in re-engagement campaigns, discovery surfaces, abandoned-cart recovery, and recommendation explanations. They are less likely to hold in low-variance transactional alerts such as password resets, delivery confirmations, payment receipts, or fraud warnings, where factual precision and latency dominate stylistic variation.

From a systems perspective, the decision to deploy LLM generation should be framed as a cost-benefit problem. Let the incremental value of LLM generation be defined as the expected gain in downstream business utility relative to a strong template or slot-filling baseline, net of additional inference, retrieval, and safety costs. In high-volume systems, even modest latency or compute overhead can dominate the marginal uplift from semantic diversity. Thus, the relevant comparison is not against manual copywriting, but against optimized templating systems that already incorporate strong ranking, personalization features, and send-time policies.

A mature architecture should therefore support dynamic routing between generation strategies. The budget allocator in our proposed pipeline (Section 9) can be interpreted not merely as a compute-control mechanism, but as a policy layer for deciding whether an LLM should be invoked at all. This reframes LLMs as one component in a portfolio of messaging strategies rather than the default endpoint for all notification composition tasks.

\section{Evaluation Frameworks Beyond Short-Term CTR}

A major weakness in the current literature is that evaluation strategies lag behind architectural complexity. Most published results still report short-term click or conversion gains, yet LLM-based notification systems operate in environments where treatment effects are delayed, interference is common, and user welfare is multi-dimensional. Figure \ref{fig:evaluation} outlines a proposed multi-layer evaluation framework.

\begin{figure}[h]
    \centering
    \includegraphics[width=\linewidth]{fig4_evaluation.png}
    \caption{Multi-Layer Evaluation Framework for LLM-Based Notification Systems.}
    \label{fig:evaluation}
\end{figure}

\subsection{The Identification Problem: Why Simple A/B Testing Is Not Enough}

Randomized A/B testing remains the operational gold standard, but its interpretation in notification systems is more fragile than commonly assumed. In networked and recommender-driven environments, the Stable Unit Treatment Value Assumption (SUTVA) is frequently violated \cite{luo2024survey} [T1]. Notifying one user can alter content popularity, inventory exposure, or social activity in ways that indirectly affect untreated users. This creates spillover effects that contaminate causal attribution. 

In addition, standard message-level tests often fail to isolate which part of the pipeline caused the uplift. A measured gain may be driven by better ranking, better timing, or better targeting rather than by language generation itself. A rigorous experimental program should therefore distinguish at least three treatment layers: item selection, send-time decision, and message realization.

\subsection{Offline Evaluation Is Structurally Biased}

Offline evaluation remains useful for rapid iteration, but it is subject to severe bias. Historical logs only contain feedback for actions that were actually taken, meaning that novel prompts or candidate framings have no unbiased counterfactual labels. This problem is especially acute for reward models trained on logged notification outcomes, because these models inherit the exposure policy that generated the data. 

Selection bias also appears at the population level. Users who permit notifications are typically more engaged and more tolerant of outreach than the marginal users whose behavior platforms most want to influence. Evaluating generation quality only on this observed subpopulation leads to optimistic results that may not transfer to churn-risk or dormant users.

\subsection{Heterogeneous Treatment Effects Must Be First-Class}

Average Treatment Effects (ATEs) are often too coarse for notification systems. Different users respond differently not only to content categories but also to persuasion intensity, timing frequency, and stylistic framing. A gain that is positive on average may conceal harm to important subgroups, including low-engagement users, new users, or users already close to opt-out thresholds. Tu et al. (2021) demonstrated that the causal effect of notifications varies drastically across user cohorts \cite{tu2021personalized} [T1]. Evaluation should therefore estimate Heterogeneous Treatment Effects (HTEs) across user segments defined by lifecycle stage, baseline engagement, CLV, and prior notification saturation.

\subsection{Multi-Objective Evaluation Is Required}

Notification quality cannot be reduced to CTR. A system optimized only for click propensity will often converge toward clickbait-like phrasing, over-alerting, or manipulative urgency cues \cite{jannach2023survey} [T1]. These behaviors may improve short-term metrics while degrading long-term trust, retention, and channel health. We recommend structuring metrics into three layers:

\begin{enumerate}
    \item \textbf{Primary outcome metrics:} CTR, CVR, order value, session return.
    \item \textbf{Guardrail metrics:} Notification opt-out rate, mute actions, dismissals without open, complaint rate, and uninstall risk.
    \item \textbf{Long-horizon health metrics:} 7-day and 30-day retention, notification fatigue trajectories, cohort survival, and reactivation durability.
\end{enumerate}

\section{Ethical Frameworks and Governance}

The ethical risks of LLM-based notification systems cannot be reduced to hallucination alone. Because notifications are interruptive, personalized, and often optimized for behavior change, they sit close to the boundary between helpful communication and covert manipulation.

\subsection{Persuasive Asymmetry and the Manipulation Boundary}

Digital platforms possess profound informational advantages over users. Susser et al. (2019) define online manipulation as the use of information technology to covertly influence decision-making by exploiting cognitive biases, bypassing rational agency \cite{susser2019technology} [T1]. When LLMs generate personalized text utilizing scarcity heuristics (``Only 1 left!'') or false urgency, they cross the boundary from persuasion into manipulation. Ethical systems must prohibit prompt instructions that encourage deceptive framing, classifying them as ``dark patterns'' \cite{mathur2021what} [T1].

\subsection{Fairness and Differential Susceptibility}

Fairness in recommender systems is often discussed in terms of exposure or ranking parity \cite{wang2023survey} [T1], but notification generation adds a different concern: differential susceptibility. Some user groups (e.g., older adults, low-literacy users) may be more influenced by certain tones, urgency cues, or authoritative language patterns than others. Fairness requires auditing reward models to ensure they do not systematically exploit the vulnerabilities of specific demographic groups to artificially inflate CTR.

\subsection{Accountability and Regulatory Compliance}

Under the EU's Digital Services Act (DSA) Article 40 and GDPR Article 22, platforms face increasing obligations regarding algorithmic transparency and the right of users not to be subject to decisions based solely on automated processing \cite{fabbri2023ethical} [T1]. Platforms must ensure data minimization in RAG contexts and provide researchers with access to audit the fairness of the notification pipeline. Frequency caps, forbidden content policies, factuality checks, and audit logs are not optional accessories; they are part of the core production design.

\section{Proposed Unified Architectural Framework}

Based on the surveyed literature and the preceding analysis, we propose an upgraded unified operational framework for LLM-based notification composition (Figure \ref{fig:pipeline}). The novelty of this framework is integrative rather than algorithmic. It specifies how known components should be sequenced and coordinated to support scalable, grounded, and policy-compliant notification delivery.

\begin{figure*}[t]
    \centering
    \includegraphics[width=\textwidth]{fig2_pipeline.png}
    \caption{The proposed unified notification pipeline, demonstrating the integration novelty of separating generation from adjacent systems and the inclusion of explicit safety guardrails, contextual bandits, and counterfactual feedback loops.}
    \label{fig:pipeline}
\end{figure*}

\subsection{Architectural Flow and Failure Modes}

The pipeline begins with three inputs: user profile/value estimates, content inventory, and contextual signals. These flow into a content-selection layer. Crucially, a \textbf{Budget-Aware Router} then decides whether the user-item-context tuple justifies invocation of the LLM pathway or whether a simpler templated strategy should be used.

If the LLM pathway is selected, \textbf{RAG} retrieves relevant context. The generator produces multiple candidates, conditioned through \textbf{PEFT/LoRA} adapters. A \textbf{Factuality and Policy Guard} stage filters candidates containing unsupported claims or prohibited phrases. The remaining candidates are scored by a \textbf{Pairwise Reward Model}. After ranking, a \textbf{Diversity and Frequency Control Layer} ensures the chosen message does not contribute to alert fatigue. Finally, a \textbf{Send-Time Optimization Module (Bandit)} determines whether the message should be delivered immediately or deferred. Outcomes then flow back into an \textbf{Online Learning Layer}.

\textbf{Failure Modes and Tradeoffs:} This pipeline trades increased inference latency and compute cost for semantic diversity. If the Factuality Guardrail is too strict, it throttles throughput; if too loose, it risks brand damage. If the Reward Model suffers from exposure bias, it will continually select the same ``safe'' candidates, defeating the purpose of the LLM generator entirely.

\section{Threats to Validity}

Several threats to the validity of this survey must be acknowledged:
\begin{enumerate}
    \item \textbf{Publication Bias:} Industry labs are highly incentivized to publish successful A/B tests and suppress negative results, potentially skewing the reported efficacy of LLM notifications.
    \item \textbf{Temporal Degradation:} The rapid evolution of LLM capabilities means that technical constraints (e.g., inference latency) reported in 2024--2025 papers may be resolved by hardware advances in 2026.
    \item \textbf{Domain Representation:} Our corpus heavily favors Social Media and E-Commerce; findings may not generalize to high-stakes domains like healthcare or financial alerts.
    \item \textbf{Reproducibility:} Many of the most important industrial systems depend on proprietary data and internal experimentation frameworks, lacking widely accepted open benchmarks.
\end{enumerate}

\section{Conclusion}

This paper surveyed the emerging field of LLM-based notification message composition and argued that its real significance lies not in language generation alone, but in the interaction between generation and the surrounding notification architecture. The literature shows that grounded generation, efficient style adaptation, relative candidate ranking, diversity control, and contextual timing can together improve the quality and relevance of notifications. 

A central conclusion of this survey is that LLM text generation is not universally the binding constraint in notification systems. For transactional, factual, or low-variance alerts, high-quality templates remain the superior design. The strongest case for LLM deployment arises when user response is sensitive to framing, when multiple context signals can be synthesized into useful copy, and when the expected business value justifies the additional inference, evaluation, and governance burden.

Future research must move beyond binary claims that LLMs do or do not improve notifications. Progress will depend on stronger experimental design, component-isolated evaluation, better off-policy learning infrastructure, more rigorous fairness auditing, and open benchmarks that allow researchers to disentangle the contribution of generation from the broader recommender stack.

\bibliographystyle{ACM-Reference-Format}
\bibliography{acmart}

\end{document}
